\section{Proof ledgers}
\label{app:proof-ledgers}

\subsection{Closure and the support quotient}

Iterating \eqref{eq:edge-recurrence} from \(e_0=0\) gives
\[
e_N-e_0=
\sum_{j=0}^{N-1}(-1)^{N-1-j}
\bigl(q+c_{\sigma(E_j)}-c_j\bigr).
\]
The coefficient of \(q\) is
\(\sum_{j=0}^{N-1}(-1)^{N-1-j}=0\) because \(N=2n\) is even.
If \(j_*\) is the unique inverse-image index of the closing edge, the two
indicator sums reduce the closure condition to
\[
(-1)^{N-1-j_*}-1=0.
\]
Thus \(j_*\) is odd.  For the support \(i\mapsto i+k\), one has
\(j_*=N-1-k\), so \(k\) is even.  For \(i\mapsto k-i\), one has
\(j_*=k\), so \(k\) is odd.  This gives \(2n\) supports and three
forced-scale lifts per support.

The phase of \(r^n\) is zero on even coordinates and one on odd
coordinates, up to a common projective scalar.  Hence
\[
r^n=\operatorname{diag}(1,\rho,1,\rho,\ldots,1,\rho),
\]
which has order three.  Together with the order-\(n\) support of \(r\),
this proves \(|r|=3n\).  The support and phase identities
\(s^2=1\) and \(srs=r^{-1}\) then give a dihedral subgroup of size
\(6n\), equal to the exhaustive list.

\subsection{Fixed-point congruences}

From \eqref{eq:galois-action},
\[
\delta(r^k)=r^{-k},\qquad
\delta(r^ks)=r^{1-k}s.
\]
The number of rotation solutions to \(2k=0\) modulo \(3n\) is
\(\gcd(2,3n)\); the reflection congruence \(2k=1\) has a solution exactly
when \(3n\) is odd.  These alternatives yield the two elements listed in
Theorem~\ref{thm:rational-form} and no others.

\subsection{Rank reduction}

The two necessary congruences
\[
n\mid4e_n,\qquad n\mid4o_n
\]
and the identity \(o_n=2(e_n-3)\) imply
\[
n\mid8e_n,\qquad n\mid8(e_n-3),\qquad n\mid24.
\]
There are therefore only seven candidate divisors \(n\ge2\).  The residue
table in the proof of Theorem~\ref{thm:denominator} is exhaustive, not a
finite approximation to an unbounded scan.

\subsection{Character reconstruction}

For any class function \(\chi\) on \(G_3\), the one-dimensional
multiplicities are
\begin{align*}
m_{\one}
&=\frac1{18}\left(\sum_{k=0}^{8}\chi(r^k)
 +\sum_{k=0}^{8}\chi(r^ks)\right),\\
m_{\varepsilon}
&=\frac1{18}\left(\sum_{k=0}^{8}\chi(r^k)
 -\sum_{k=0}^{8}\chi(r^ks)\right).
\end{align*}
The two-dimensional multiplicities are
\[
m_{U_j}=\frac1{18}\sum_{k=0}^{8}
\chi(r^k)\bigl(\zeta_9^{-jk}+\zeta_9^{jk}\bigr),
\qquad 1\le j\le4.
\]
Applying these formulas to the two trace rows gives
\begin{center}
\begin{tabular}{@{}c|rrrrrr@{}}
\toprule
character & \(m_{\one}\) & \(m_{\varepsilon}\) & \(m_{U_1}\) &
\(m_{U_2}\) & \(m_{U_3}\) & \(m_{U_4}\)\\
\midrule
\(H^{2,1}(X_3)\) & \(0\) & \(2\) & \(2\) & \(2\) & \(3\) & \(2\)\\
\(\Epack_3\) & \(1\) & \(2\) & \(3\) & \(3\) & \(1\) & \(3\)\\
\bottomrule
\end{tabular}
\end{center}
Dimension checks give \(20\) and \(23\), respectively.  The Galois action
on ninth roots permutes \(U_1,U_2,U_4\) and leaves the multiplicity pair
\((3,4)\) after the odd rail is doubled.  This is why splitting that orbit
into three unrelated rational sectors is not allowed.
